% Modernised reading — fonds Alexandre Grothendieck, shelfmark 22, pages 1-29.
%
% This is an INTERPRETATION, not a transcription. It re-reads the pages in
% today's notation and vocabulary, states the hypotheses the manuscript leaves
% implicit, and drops the critical apparatus so that the mathematics reads
% straight through. Everything the brackets and underlines carried moves into a
% footnote. It is held to being mathematically correct as it stands; where the
% manuscript is loose or mistaken, the true statement is what appears here and
% the footnote says what the page has. In any dispute of substance the
% transcription governs: cite batch-01/02.fr.tex.
%
% Pass: Opus 5 (claude-opus-5), 2026-09-19 — first pass, unchecked against the
% pages by a human. The model was chosen explicitly for this pass.
%
% Scope: one document for the shelfmark, its two batches read as one argument.
% Twenty-two of the twenty-nine leaves carry something; pages 4, 6, 7, 8 are
% blank versos of the offprints, 18 and 25 are covers carrying one word each,
% and 27 and 29 are printed matter foreign to the folder (they are the versos
% of the pencil leaves 26 and 28).
%
% Spine. The folder asks ONE question in four registers, and the question is:
% WHEN IS A LOCAL RING SMOOTH OVER ITS BASE IN THE PURELY FORMAL SENSE — every
% solution modulo a square-zero ideal lifts — AND WHAT MEASURES THE
% OBSTRUCTION? The answer the folder converges on is 3.13: formally smooth =
% flat + geometrically regular closed fibre. The obstruction, in every register,
% is a module of differentials.
%
%   pp. 1-8    TWO PRINTED NOTES OF J.-P. JOUANOLOU (CRAS 1963), not
%              transcribed; only his marginalia, on pages 1, 2, 3 and 5. Both
%              notes are about the DIFFERENTIAL CRITERION OF REGULARITY (Nakai,
%              Suzuki, Mount, Kunz: D(A) free iff A regular) and about
%              D_A(B) = 0. That is the published state of the art on exactly
%              the fibre condition 3.13 will need. His annotations are a
%              triage: « connu », « connu (Nagata) », « faux » twice,
%              « C ⊃ A !!! », « séries formelles ?? », « assez rare ! »,
%              « non vérifié ».
%   pp. 9-17   THE TYPESCRIPT, his pages 24 to 32, the numéro 3 of a chapter,
%              addressed to a « tu ». Lemme 3.2 (criterion for a continuous
%              homomorphism to be a topological isomorphism), Lemme 3.3 and
%              Cor. 3.4, 3.5 (products of flats are flat; a strict countable
%              limit of projectives is flat), Théorème 3.6 (formal smoothness
%              and uniqueness of the lift), Déf. 3.7 (Cohen algebra), Cor. 3.8
%              (uniqueness + four equivalent forms), Cor. 3.9 (the local
%              criterion), Cor. 3.10 (EXISTENCE, via Cohen's structure
%              theorem), Cor. 3.11, 3.12 (the unramified case), Cor. 3.13 (the
%              punchline), and a page of self-criticism proposing to throw the
%              whole numéro away and keep 3.13 alone.
%   pp. 18-20  A cover inscribed « Questions », then four OPEN QUESTIONS in
%              ink. Question 2) is the one pages 21-24 answer.
%   pp. 21-24  Ink. IS A = K[[t_1,…,t_n]] FORMALLY SMOOTH OVER k? Case p > 0
%              settled by p-bases; case p = 0, « Je suis pessimiste », never
%              for the discrete topology, with the obstruction exhibited over
%              k = Q((x_i)).
%   pp. 25-28  A cover inscribed « Différentielles » (forward-referenced from
%              p. 21 as « cf. "Diff" »), then pencil, a later campaign: the
%              square k → K, k ∩ K^p → K^p, the injectivity of
%              Ω^1_{k/k∩K^p} ⊗_k K → Ω^1_K, and a counterexample in
%              characteristic p.
%
% What only a folder-wide reading can say. Three links, none stated on any one
% leaf:
%   - QUESTION 2) OF PAGE 19 IS ANSWERED ON PAGES 21 TO 24. The wording is
%     nearly verbatim ("est-il ft simple pour sa top. discrète ?" against "A
%     est-il … pour la topologie discrète ???"), and the answer is: no in
%     characteristic zero as soon as dim A > 0, yes in characteristic p when
%     K has a finite p-basis. This is the single strongest piece of evidence
%     that the folder is one dossier and not four.
%   - THE DIVIDING LINE IS THE DIMENSION OF Ω_{K/k}. Page 23's proposition
%     assumes it finite; page 21's marginal notes Ω_{K/k} free; page 22's
%     counterexample makes it infinite. The three campaigns agree on one
%     invariant without ever saying so.
%   - THE OFFPRINTS ARE THE FIBRE CONDITION. 3.13 needs "B ⊗_A k universally
%     regular over k"; Jouanolou's two notes are about when regularity is
%     detected by differentials. Nothing in pp. 9-28 cites Jouanolou, so this
%     link is thematic only, and is stated as such in the body.
%
% NOTATION fixed for the folder.
%   - « formellement simple » is his word; it is written FORMELLEMENT LISSE
%     throughout, which is today's name and which the folder itself drifts to
%     on p. 22. The rename is footnoted once, at its first occurrence.
%   - r(A), his notation for the radical of a local ring, is written 𝔪_A.
%   - « universellement régulier » is kept as his term and glossed once as
%     geometrically regular.
%   - « condition P » is NOT defined in this folder (it belongs to an earlier
%     numéro). It is kept under his name, and exactly what the folder uses of
%     it is isolated in the spine section.
%   - « multiplicité radicielle finie » is likewise not defined here. Kept as
%     his term; what it is FOR is stated (the NB of p. 14 says it: to know B_0
%     is Cohen-Macaulay).
%   - A_i = A/I_i, B_i = B/I_iB, C_i = C/I_iC, B_0 = B/IB throughout.
%
% Substantive departures, each footnoted where it occurs:
%   - p. 11 (Théorème 3.6 (ii)): the five equivalent conditions are typed about
%     B and must be about C. As typed they are consequences of the theorem's
%     own hypotheses and the equivalence is vacuous; the proof, which lifts
%     u_0^{-1} : C_0 → B_0 using formal smoothness of C, settles it. Written
%     here about C.
%   - p. 11: the proof says assertion (i) follows "de a)"; it needs a) AND the
%     handwritten clause b), which was inserted after the proof was typed.
%   - p. 9: the marginal calls hypothesis d) useless, but the typed proof of
%     (ii) ⟹ (i bis) uses it. The hypothesis is kept.
%   - p. 10 (Lemme 3.3): the handwritten « de type fini » on the free modules
%     is not an embellishment, it is what makes the proof work. Kept, and the
%     necessity is stated.
%   - p. 15 (Cor. 3.10): « le théorème de structure 3.4 » cannot be the 3.4 of
%     this numéro (products of flat modules). It is Cohen's structure theorem,
%     from numéro 2.
%   - p. 17 (Cor. 3.13): the handwritten alternative hypothesis is written
%     "k(B) = k(A)B"; page 15 writes the same clause "r(B) = r(A)B", which is
%     the hypothesis of 3.12 and the only reading that makes it an alternative.
%   - p. 22: « de dim A ≥ 0 » is read as dim A > 0. Dimension zero must be
%     excluded: in characteristic zero every field extension is formally
%     smooth.
%   - p. 23: the right half of the Proposition is illegible. The conclusion
%     given here is the one its hypothesis and the surrounding argument force,
%     and it is OURS.
%   - p. 24: the reduction to D(t_i) = 0 is carried out on pp. 22 and 23 in
%     prose that is largely illegible. Only its shape is given.
%
% NUMBERING. The digit opening each statement is typed then gone over in ink
% and reads as 3 or 4. The transcription retained 3.N. That reading is taken
% up here and strengthened; the argument is in a footnote in the spine section.
%
% What the folder announces and never establishes: the definition of condition
% P; the definition of « multiplicité radicielle »; the statements 3.1, 2.3,
% 2.6, 2.12 and 0_III 10.3.1, all cited and none present; the head of Lemme 3.2
% (his page 23, absent); the four open questions of pp. 19-20, whose connective
% prose and whose two long marginal additions could not be read at all; the end
% of the counterexample of p. 24, which stops on the construction; and the role
% of the norm N_{K/k}, named twice on p. 26 in illegible surroundings.
\documentclass[11pt,a4paper]{article}
% Le préambule commun à toutes les transcriptions du fonds Grothendieck.
%
% Il tient en une idée : **l'incertitude se compose, elle ne se résout pas.**
% Une transcription qui lisse un mot douteux a détruit la seule chose pour
% laquelle on la faisait — un lecteur doit pouvoir savoir, à chaque endroit,
% ce qui était sur le feuillet et ce qui a été ajouté. Les six macros qui
% suivent sont donc l'appareil critique, et non de la décoration.
%
% Les mêmes commandes sont reconnues par `scripts/render.mjs`, qui produit la
% vue de lecture du site. Une macro ajoutée ici doit l'être aussi là-bas,
% faute de quoi le rendu échouera bruyamment — ce qui est voulu : un rendu
% silencieusement faux est pire qu'un rendu absent.

\ProvidesPackage{grothendieck}[2026/08/08 Transcriptions du fonds Grothendieck]

\usepackage{amsmath,amssymb,amsthm}
\usepackage{mathrsfs}
\usepackage{stmaryrd}
% Les diagrammes commutatifs sont partout dans ce fonds : c'est la langue
% même de la géométrie algébrique de Grothendieck, et une transcription qui
% les rendrait en prose ne transcrirait rien.
\usepackage{tikz-cd}
% Français par défaut — c'est la langue des feuillets — mais l'anglais est
% chargé aussi : la traduction et le résumé sont des documents anglais, et la
% typographie française (espaces avant les deux-points) y serait une faute.
% Ces documents basculent par \selectlanguage{english} en tête de corps.
\usepackage[english,french]{babel}
\usepackage{fontspec}
\usepackage{geometry}
\geometry{a4paper,margin=25mm,marginparwidth=28mm}
\usepackage{marginnote}
\usepackage{xcolor}
% ulem, not soul. `soul`'s \st and \ul are built on a character-by-character
% scan that XeLaTeX's font machinery breaks: the document fails with an
% undefined control sequence rather than degrading. ulem does the same two jobs
% and is Unicode-engine safe. `normalem` stops it hijacking \emph, which the
% transcriptions use for Grothendieck's own emphasis.
\usepackage[normalem]{ulem}
\usepackage{hyperref}
\hypersetup{colorlinks=true,linkcolor=black,urlcolor=black,pdfborder={0 0 0}}

% --- Métadonnées du lot -----------------------------------------------------
% Reprises telles quelles par le script de rendu, qui les affiche en tête de
% la vue de lecture. Elles fixent ce que le fichier prétend couvrir : sans
% elles, une transcription flotte sans point d'attache dans un fonds de seize
% mille feuillets.

\newcommand{\folder}[1]{\gdef\grothFolder{#1}}
\newcommand{\batch}[1]{\gdef\grothBatch{#1}}
\newcommand{\leaves}[2]{\gdef\grothFirst{#1}\gdef\grothLast{#2}}
\newcommand{\foldertitle}[1]{\gdef\grothFoldertitle{#1}}
\gdef\grothFolder{?}\gdef\grothBatch{?}\gdef\grothFirst{?}\gdef\grothLast{?}\gdef\grothFoldertitle{}

% --- Le résumé --------------------------------------------------------------
%
% Il ouvre la lecture modernisée, sous le titre « Résumé », et il est composé
% en petit. Ce n'est pas un ornement : c'est le seul passage du document qui
% ne soit pas des mathématiques, et un lecteur doit pouvoir le reconnaître
% avant de l'avoir lu — puis le sauter s'il connaît déjà le sujet, ou s'y
% arrêter s'il ne le connaît pas. Le filet qui le suit marque la reprise du
% corps.

\newenvironment{resume}
  {\small\setlength{\parskip}{0.45em}}
  {\par\medskip\hrule\medskip}

% --- Mots-clés ---------------------------------------------------------------
%
% En anglais, en clôture du résumé de la lecture modernisée : le vocabulaire
% moderne sous lequel chercher ce que les feuillets construisent. Le manifeste
% du site (`npm run manifest`) les extrait pour étiqueter la cote entière —
% cette ligne est la source unique des tags, il n'y a pas de fichier à côté.
\newcommand{\keywords}[1]{\par\smallskip\noindent
  {\footnotesize\textsc{Keywords} --- \textit{#1}}\par}

% --- L'appareil critique ----------------------------------------------------

% Le numéro de feuillet, en marge. C'est l'ancre qui permet de régler un
% désaccord sur une formule en regardant *un* feuillet plutôt qu'une chemise
% de six cents. Le site s'en sert aussi pour tourner les pages du fac-similé
% quand on fait défiler la transcription.
\newcommand{\leaf}[1]{%
  \marginnote{\footnotesize\color{gray!70}\textsf{#1}}%
  \label{leaf:#1}\ignorespaces}

% Illisible : on ne devine pas. Un mot inventé qui se lit comme les autres est
% le pire résultat possible.
\newcommand{\ill}{\textcolor{red!60!black}{[\,\ldots\,]}}

% Lecture douteuse : le mot est donné, et signalé comme incertain.
\newcommand{\uncertain}[1]{\uline{#1}}

% Ajout de l'éditeur — une lettre manquante, un mot sous-entendu.
\newcommand{\add}[1]{\textcolor{blue!50!black}{[#1]}}

% Biffé par Grothendieck lui-même. Ce qu'il a rayé fait partie du document :
% une rature dit souvent où la pensée a changé de direction.
\newcommand{\struck}[1]{\sout{#1}}

% Note du transcripteur, sur l'état matériel du feuillet.
\newcommand{\note}[1]{\par\smallskip{\small\color{gray!60!black}\textit{[#1]}}\par\smallskip}

% Note marginale de Grothendieck, distincte du corps du texte.
\newcommand{\marginal}[1]{\par\smallskip\noindent\hspace{1em}%
  {\small\color{gray!60!black}#1}\par\smallskip}

% --- Titre ------------------------------------------------------------------

\AtBeginDocument{%
  \begin{center}
    {\large\bfseries Fonds Alexandre Grothendieck --- cote n\textsuperscript{o}~\grothFolder}\\[2pt]
    {\itshape\grothFoldertitle}\\[4pt]
    {\small feuillets \grothFirst--\grothLast\ (lot \grothBatch)}\\[2pt]
    {\footnotesize\color{gray!60!black}Transcription établie d'après le fac-similé numérisé
     par l'Université de Montpellier.\\
     Les lectures incertaines sont soulignées, les illisibles notées [\,\ldots\,],
     les ajouts entre crochets.}
  \end{center}
  \vspace{1em}}

\endinput


\folder{22}
\pages{1}{29}
\dating{1963-[vers 1967]}
\watermark{Édition de démonstration\\interprétation personnelle de l'œuvre}
\foldertitle{[EGA IV]. [Chapitre] 0 IV. Cohen, différentielles, algèbres formellement simples (Compléments) : tirés à part (1963), tapuscrit annoté (s.d.), notes manuscrites (s.d.). — lecture modernisée du dossier entier}

\begin{document}

\section*{Résumé}

\begin{resume}
Il y a une question qu'on peut poser d'une machine à calculer approximativement,
et c'est celle de tout ce dossier : \emph{si j'ai résolu un système d'équations
à une certaine précision, puis-je toujours corriger ma solution pour gagner un
cran de précision de plus ?}

Prenons le cas le plus familier. On cherche une racine d'un polynôme, on en
connaît une valeur approchée, et la méthode de Newton la raffine : on remplace
la valeur approchée par une valeur meilleure, et on recommence. La méthode
marche à une condition, celle que la dérivée ne s'annule pas au point considéré
— c'est-à-dire que la courbe y soit \emph{lisse}. Là où la dérivée s'annule, la
correction se bloque : il existe des solutions approchées qu'aucune correction
ne prolonge.

L'algèbre a transposé cette situation mot pour mot, et sans aucune limite ni
aucune convergence. On se donne un anneau, on se donne un idéal dont le carré
est nul — c'est la façon algébrique de dire « à un infiniment petit près » — et
l'on demande que toute solution modulo cet idéal se relève en une solution
exacte. Une algèbre pour laquelle ce relèvement est toujours possible est dite
\emph{formellement lisse}. Grothendieck, à la date de ces feuillets, l'appelle
encore \emph{formellement simple}, et l'on verra le mot changer sous nos yeux
d'un feuillet à l'autre du dossier.

Le problème auquel le dossier s'attaque est qu'on ne sait pas reconnaître les
algèbres formellement lisses à l'œil. Le relèvement doit être possible pour
\emph{tous} les anneaux et \emph{tous} les idéaux de carré nul : c'est une
infinité de vérifications, dont aucune n'est calculable. Il faut donc un
critère, c'est-à-dire une propriété qui se lise sur l'algèbre elle-même. Le
dossier en trouve un, et c'est son point d'arrivée : une algèbre locale est
formellement lisse sur sa base exactement quand elle est \emph{plate} sur elle
— la base ne se déforme pas sous ses pieds — et quand sa \emph{fibre spéciale},
ce qui reste de l'algèbre quand on annule tout ce qui vient de la base, est
\emph{régulière}, et le reste après n'importe quelle extension du corps de
base. Deux conditions, l'une de famille, l'autre de fibre, et l'on peut les
vérifier.

Une deuxième histoire traverse le dossier, celle de Cohen. Un théorème de
structure, dû à Cohen, dit qu'un anneau local complet n'a pas de mystère : il
ressemble à un anneau de séries formelles à coefficients dans un corps. La
question relative se pose aussitôt : on se donne une base, on se donne la fibre
spéciale, existe-t-il au-dessus de la base une algèbre qui redonne cette fibre,
et est-elle unique ? Le tapuscrit du dossier appelle \emph{algèbre de Cohen} la
bonne notion pour poser cette question, et il y répond en deux temps :
l'unicité est acquise sans rien supposer de plus, l'existence ne l'est que si
l'extension de corps résiduels en jeu est séparable. Le cas d'école est celui
où la fibre spéciale se réduit à une extension de corps : on retrouve alors
l'énoncé qui gouverne les anneaux de Cohen et, à travers eux, la façon dont
tout anneau local complet se laisse écrire.

La troisième histoire est celle des \emph{différentielles}, et c'est elle qui
donne au dossier sa surprise. Une différentielle est l'objet algébrique qui
enregistre les directions dans lesquelles on peut dériver. Or il y a deux
manières de demander qu'une algèbre soit formellement lisse, selon qu'on
s'autorise ou non à passer à la limite : la topologie \emph{habituelle}
d'un anneau local permet de dire que deux éléments sont proches quand leur
différence est très divisible ; la topologie \emph{discrète} l'interdit et
exige un relèvement exact, en une fois. Grothendieck se demande si la première
entraîne la seconde. Quatre feuillets manuscrits répondent : non. La raison est
d'une simplicité désarmante et se voit sur un exemple. Prenons pour corps de
base un corps engendré par une infinité de variables indépendantes, et
au-dessus de lui l'anneau des séries formelles en une variable. Dériver une
série formelle, c'est dériver chacun de ses coefficients ; si chaque
coefficient fait intervenir une variable neuve, la dérivée de la série fait
intervenir une infinité de directions à la fois. Or l'objet où la dérivée doit
atterrir n'autorise que des \emph{sommes finies} de directions. La dérivée
existe donc terme à terme, elle converge dans un espace un peu plus grand, et
elle n'atterrit nulle part dans celui qu'on lui impose. L'obstruction n'est pas
un défaut de la géométrie, c'est un défaut de place.

En caractéristique positive, la même question reçoit la réponse inverse, et
pour une raison également simple : le corps de base y possède, quand il n'est
pas trop gros, un système fini de directions indépendantes — une
\emph{p-base} — et une infinité de directions ne peut pas se produire. C'est
la même mesure qui décide dans les deux cas : la dimension de l'espace des
différentielles du corps. Finie, tout va bien ; infinie, rien ne va.

Reste à dire ce que ces feuillets sont, car cela fait partie de ce qu'on y lit.
Le dossier est un dossier de travail, et il montre quatre régimes d'écriture
que rien n'oblige à cohabiter. Deux tirés à part imprimés, dont le texte n'est
pas de lui, et qu'il annote en juge : « connu », « connu (Nagata) », « faux »,
et en grand, en haut d'une page, « $C \supset A$ !!! ». Un tapuscrit envoyé à
quelqu'un qu'il tutoie, où il écrit un chapitre et plaide simultanément contre
lui : « je me repens un peu de la définition 3.7 », « la terminologie “algèbre
de Cohen” risque plutôt d'embrouiller les idées (d'autant plus que nous
n'utiliserons plus jamais que “formellement simple”) », et pour finir le
conseil de tout jeter sauf un énoncé, de le démontrer « en deux pages ou
trois », « et un point c'est tout ». Une note en bas de feuillet où Serre lui
apprend qu'un article de Mac Lane de 1940 rend inutile une hypothèse qu'il
vient de faire, avec la commission passée au destinataire d'aller y voir. Et
quatre feuillets d'encre où il écrit, avant de savoir : « Je suis
pessimiste ». C'est le moment d'une pensée mathématique où le pressentiment
précède la démonstration, et il est rare qu'un tel moment soit conservé par
écrit.

Les noms sous lesquels on cherchera la suite : lissité formelle et morphismes
réguliers, théorème de structure de Cohen et anneaux de Cohen, critère local de
platitude, p-bases et théorème de Cartier, anneaux excellents, ouverture du
lieu régulier, et — puisque le dernier feuillet de questions s'achève sur ces
mots soulignés — la résolution des singularités.

\keywords{formally smooth morphism, formal smoothness, formally unramified, formally etale, Cohen structure theorem, Cohen ring, adic ring, ideal of definition, preadmissible ring, completed tensor product, local criterion of flatness, flat module, faithfully flat descent, finitely presented module, product of flat modules, strict projective system, projective limit, regular local ring, regular system of parameters, Cohen-Macaulay ring, geometrically regular fibre, regular homomorphism, module of Kahler differentials, differential criterion of regularity, p-basis, p-independence, imperfect field, separable field extension, inseparable extension, radicial extension, first fundamental exact sequence, deformation and lifting, excellent ring, openness of the regular locus, resolution of singularities, analytically reduced local ring, integral closure finiteness}
\end{resume}

\pagerange{1}{29}
\section*{Le fil du dossier, et les conventions}

\emph{Les quatre campagnes.} Le dossier n'est pas un texte, c'est un dossier :
quatre campagnes d'écriture de dates et de natures différentes, rangées dans
une même chemise.

\begin{enumerate}
  \item \emph{Pages 1 à 8}, imprimé. Deux notes aux Comptes rendus de
  J.-P. Jouanolou, celle du 4 novembre 1963 (t. 257, p. 2948-2950) et celle du
  14 janvier 1963 (t. 256, p. 842-844). Le texte imprimé n'est pas reproduit
  dans la transcription ; seules figurent ses annotations, sur les pages 1, 2,
  3 et 5.
  \item \emph{Pages 9 à 17}, tapuscrit paginé par lui 24 à 32, avec une
  campagne de corrections à l'encre. C'est le numéro 3 d'un chapitre, adressé à
  un correspondant qu'il tutoie et que le dossier ne nomme pas. Du lemme 3.2
  au corollaire 3.13.
  \item \emph{Pages 18 à 20}, encre. Une chemise portant le seul mot
  « Questions », puis quatre questions ouvertes.
  \item \emph{Pages 21 à 28}, en deux temps : l'encre des pages 21 à 24, sur
  le verso de feuilles à en-tête de sa dactylographe ; puis, après une chemise
  portant le seul mot « Différentielles » (page 25), le crayon des pages 26 et
  28, d'une campagne visiblement postérieure.
\end{enumerate}

Les pages 4, 6, 7 et 8 sont les versos blancs des tirés à part. Les pages 27 et
29 portent un imprimé sur la topologie de l'espace des répartitions et le
groupe des idèles, étranger au chapitre et sans annotation de lui : ce sont les
versos des feuillets 26 et 28, et ils ne sont pas lus ici.

\emph{Le dossier tient-il ensemble ?} Le titre d'inventaire annonce trois
choses, « Cohen, différentielles, algèbres formellement simples », et il faut
dire honnêtement lesquelles des quatre campagnes se répondent et lesquelles ne
font que voisiner.

Deux liens sont documentés par les feuillets eux-mêmes, et ils sont solides. Le
premier : la question 2) de la page 19 demande si une algèbre locale
noethérienne formellement lisse sur $k$ pour sa topologie habituelle l'est
encore pour la topologie discrète ; les pages 21 à 24 posent exactement cette
question dès leur première ligne et y répondent. Le second : la page 21 renvoie
explicitement, entre crochets, à « Diff », c'est-à-dire à la chemise
« Différentielles » qui ouvre la campagne au crayon quatre feuillets plus loin.
Les pages 19 à 28 forment donc bien, malgré la différence d'encre et de date,
une seule enquête, et son unité est écrite sur les feuillets.

Le lien avec le tapuscrit est plus lâche, mais il est réel et il est
mathématique. Le tapuscrit aboutit au corollaire 3.13, dont la seconde
condition est que la fibre spéciale $B \otimes_A k$ soit universellement
régulière sur $k$ ; les questions des pages 19 et 20 portent toutes sur cette
même condition, la question 1) demandant précisément si l'on peut se dispenser
de l'hypothèse de finitude que 3.13 fait sur l'extension résiduelle. Rien dans
le tapuscrit ne renvoie aux questions, ni les questions au tapuscrit, mais elles
tournent autour du même énoncé et de ses hypothèses superflues.

Le lien avec les tirés à part, en revanche, n'est \emph{pas} documenté : aucune
ligne des pages 9 à 28 ne cite Jouanolou. Ce qu'on peut dire est thématique, et
il faut le dire comme tel. Les deux notes annotées portent l'une et l'autre sur
le critère différentiel de régularité — la régularité d'une localité se lit sur
la liberté de son module de différentielles absolu — et sur l'annulation du
module de différentielles relatif ; c'est-à-dire sur les deux ingrédients dont
3.13 a besoin, la régularité de la fibre et la non-ramification. Que ces deux
tirés à part de 1963 dorment dans la même chemise que le tapuscrit n'est pas un
hasard de rangement ; mais l'affirmer plus fortement serait fabriquer une unité
que les feuillets ne portent pas.

La conclusion est donc mesurée : le dossier tient ensemble comme dossier de
travail autour d'une question unique, la lissité formelle et ce qui la mesure ;
il ne tient pas ensemble comme texte. Les trois mots du titre d'inventaire sont
trois faces de cette question — « Cohen » est le théorème d'existence et
d'unicité, « algèbres formellement simples » la notion, « différentielles »
l'obstruction — et c'est la meilleure description qu'on puisse en donner sans
en inventer davantage.

\emph{La numérotation du tapuscrit.} Les énoncés sont numérotés $3.N$ dans tout
ce qui suit ; le chiffre de tête est, sur le film, aussi lisible $3$ que
$4$.\footnote{Le transcripteur a retenu $3.N$ parce que les renvois
dactylographiés restés intacts dans le corps de la prose donnent $3.1$, $3.5$,
$3.9$, ainsi que $2.6$ et $2.3$ pour le numéro précédent. Ce raisonnement est
repris ici, et il se renforce de trois observations que la transcription ne
tire pas. D'abord, $3.5$ est un énoncé de la série elle-même, page 10 ; le
renvoi dactylographié intact « en vertu de 3.5. » de la page 11 vise donc une
étiquette de la série, et étiquette et renvoi s'accordent sur le même chiffre.
Ensuite, la marginale manuscrite de la page 17 écrit « garder cependant 3.1. au
Nº3 » au moment même où le corps du texte parle du « présent numéro » : la
main, et non plus la machine, nomme le numéro et un de ses énoncés, et les
accorde. Enfin, puisque la transcription note que les corrections manuscrites
donnent elles aussi $3$, la surfrappe n'est pas une renumérotation du numéro :
l'indécision porte sur le dessin du chiffre, non sur sa valeur. Reste une
réserve, qu'il faut énoncer : si toute la série avait été renumérotée, les
renvois dactylographiés porteraient l'ancienne numérotation et les étiquettes à
l'encre la nouvelle, et l'argument tomberait. C'est précisément ce que l'accord
des deux campagnes exclut. Un dernier indice, moins net, va dans le même sens :
la page 15 renvoie au « théorème de structure 3.4 », qui ne peut être le 3.4 de
la série et doit être $2.4$ — une faute de frappe d'un $2$ en $3$, dans un
numéro où le $3$ était sous les doigts.}

\emph{Les noms et les notations, valables pour tout le dossier.} Ce qu'il
appelle \emph{formellement simple} est écrit ici \emph{formellement lisse},
qui est le nom d'aujourd'hui ; le dossier fait lui-même le glissement, puisque
la page 22 écrit « formellement lisse » là où le tapuscrit écrivait
« formellement simple ».\footnote{« Simple » était, dans la terminologie des
premiers \emph{Éléments}, le mot pour ce qu'on appelle aujourd'hui
\emph{lisse} ; le changement de vocabulaire est postérieur à la rédaction du
tapuscrit et antérieur aux feuillets manuscrits, ce qui donne au dossier un
repère de datation interne, indépendant de l'inventaire.} Son $r(A)$, le radical
d'un anneau local, est écrit $\mathfrak{m}_A$. Ses $A_i = A/I_iA$ sont écrits
$A_i = A/I_i$, et $B_i = B/I_iB$, $C_i = C/I_iC$, $B_0 = B/IB$ de même.
L'expression \emph{universellement régulière sur $k$} est gardée telle quelle ;
elle désigne ce qu'on appelle aujourd'hui une algèbre \emph{géométriquement
régulière} sur $k$, c'est-à-dire régulière et qui le reste après toute
extension du corps de base.

Deux notions du dossier ne sont pas définies dans le dossier, et il vaut mieux
le dire d'emblée que de le laisser découvrir.

La \emph{condition P} vient d'un numéro antérieur. Ce que le tapuscrit en
utilise tient en deux points, et c'est tout ce qu'on en saura ici : elle
implique la platitude, par le corollaire 3.5, dès que l'anneau de base est
noethérien ; et elle est impliquée par la platitude sous les hypothèses
noethériennes de 3.8, dont elle devient alors une reformulation. Le passage
biffé du corollaire 3.5 laisse voir ce qu'elle disait : un module topologique
$P$ satisfait P lorsqu'il est isomorphe à la limite projective d'un système
projectif \emph{strict}, admettant un système cofinal dénombrable, de modules
projectifs de type fini.\footnote{Le passage est frappé de x sur deux lignes et
demie et reste partiellement lisible ; on y lit « de $A$-modules projectifs de
type fini », et les mots « stricte » et « dénombrable » y sont glosés. La
définition donnée ici est donc une reconstitution à partir d'un passage
supprimé, et elle n'est pas nécessaire à la suite : les énoncés du numéro
n'utilisent de P que les deux implications dites.}

La \emph{multiplicité radicielle finie} d'une extension résiduelle n'est pas
définie non plus. Ce qu'on sait d'elle par le dossier : elle est vérifiée par
les extensions séparables et par les extensions de type fini (3.9 et 3.13 le
disent) ; elle ne sert qu'à garantir que la fibre spéciale est de
Cohen-Macaulay (le NB de la page 14 le dit explicitement) ; et il la croit
superflue.

\pagerange{1}{8}
\section*{I. Les deux notes annotées de Jouanolou (pages 1 à 8)}

Huit feuillets imprimés ouvrent le dossier, et leur texte n'est pas de lui. Ce
sont deux notes de J.-P. Jouanolou aux Comptes rendus, toutes deux de 1963 :
« Étude différentielle des anneaux locaux réguliers : applications », présentée
le 4 novembre, et « Sur les modules de différentielles », présentée le 14
janvier. Elles appartiennent au dossier par leur sujet et par ce qu'il en a
fait : des verdicts en marge, et rien d'autre.

Le sujet des deux notes est le \emph{critère différentiel de régularité}. Pour
une localité $A$, on dispose du module des différentielles absolu $D(A)$, et
une famille de théorèmes — Nakai, Suzuki, Mount, Kunz — relie la régularité de
$A$ à des propriétés de liberté ou de finitude de ce module. C'est très
exactement l'ingrédient dont le tapuscrit aura besoin sept feuillets plus loin,
lorsque le corollaire 3.13 demandera que la fibre spéciale soit régulière ; et
c'est la seule raison pour laquelle ces huit feuillets sont ici.

Ses annotations forment un tri, et il n'y en a pas d'autre nature. Contre le
lemme sur le noyau de $\hat{A} \otimes_A D(A) \to D(\hat{A})$ et tout le bloc
qui le suit jusqu'à la proposition 6 — les énoncés de séparabilité du corps des
fractions du complété sur celui de $A$, et le passage par le théorème de
structure de Cohen — un long trait vertical et un mot : « connu ». Contre la
proposition 7, la finitude de la fermeture intégrale de $A$ dans une extension
finie de son corps des fractions : « connu (Nagata) ». Contre le dernier alinéa
de la page 3, celui qui tire du lemme 2 de Zariski-Samuel qu'une localité
analytique normale est analytiquement normale : « connu » encore. Contre le
corollaire 3, dont l'hypothèse est que le séparé $D'(A)$ soit de type fini :
« assez rare ! ». Contre la proposition 2, sur le lieu singulier d'une algèbre
affine : « séries formelles ?? », c'est-à-dire la question de savoir si le
passage des séries convergentes aux séries formelles est licite.

Trois annotations sont des objections. En haut de la page 3, en grand, au bout
d'un trait qui descend vers la première ligne imprimée : « $C \supset A$ !!! ».
L'énoncé visé parle de « toute $A$-algèbre intègre qui est localisée d'une
$A$-algèbre finie $C$ », et l'objection est que cet énoncé suppose tacitement
$A$ contenu dans $C$. Sur la même page, un long trait horizontal souligne la
parenthèse « (Pour appliquer le lemme, on remarque que $h_1$ et $h_2$ sont
stables par localisation.) », et la marge porte : « faux ». Enfin, contre le
lemme du numéro 3 de la première note, une annotation oblique de quatre lignes
brèves qui commence par « faux », fait figurer une flèche $A \to B$ et s'achève
sur deux points d'interrogation.\footnote{Cette annotation n'est lisible qu'en
partie, et ce qui manque est ce qui expliquerait l'objection. Il en va de même
de trois autres marginales de ces feuillets, qui n'ont pas pu être lues du
tout : en marge de droite de la page 1, trois lignes courtes contre le
corollaire 2 [Kunz], dont le premier mot pourrait être « contenu » ; un mot,
contre la proposition 3, qui se termine par « -ique » ; et le premier trait
d'une annotation dont la suite se lit « hypothèse ». Aucune reconstruction n'en
est proposée ici.}

Il faut ajouter une réserve, qui vaut pour toute cette section : le texte
imprimé n'ayant pas été transcrit, aucun de ces verdicts ne peut être vérifié
sur pièces. Ce qui est établi est qu'il les a portés, non qu'ils sont fondés.

La seconde note, « Sur les modules de différentielles », ne reçoit que deux
marques, toutes deux sur son premier feuillet. Contre la définition du numéro 1
— $B$ est dite \emph{quasi-séparable} sur $A$ lorsque $D_A(B) = 0$, ce qu'on
appelle aujourd'hui \emph{formellement non ramifiée} — et contre l'affirmation
que cette propriété est locale : « non vérifié », souligné. Et en bas de page,
contre la ligne qui traduit la relation $\mathfrak{n} = \mathfrak{m}B$ par « ce
qui exprime que $B$ n'est pas ramifié sur $A$ » : un point d'interrogation.
Cette relation $\mathfrak{n} = \mathfrak{m}B$ est celle-là même qui sera
l'hypothèse du corollaire 3.12, quinze feuillets plus loin, sous la forme
$\mathfrak{m}_B = \mathfrak{m}_A B$.

\pagerange{9}{17}
\section*{II. Le tapuscrit : le numéro 3 (pages 9 à 17)}

Neuf feuillets dactylographiés, paginés par lui 24 à 32, forment une suite
continue et incomplète : elle commence au milieu d'une phrase, sa page 23
n'étant pas dans le dossier. C'est un brouillon envoyé — il y dit « j'espère
que tu pourras la rendre plus lisible », « que tu regardes cela » — et les
annotations à l'encre sont ses corrections au brouillon, portées avant ou après
l'envoi.

L'architecture du numéro est celle-ci. Un critère pour reconnaître qu'un
homomorphisme continu est un isomorphisme (3.2) ; trois énoncés de platitude
(3.3 à 3.5) qui servent à transformer une condition de présentation en une
condition de platitude ; le théorème central (3.6), qui combine les deux ; une
définition (3.7) qui nomme la situation ; puis une cascade de corollaires
(3.8 à 3.13) qui descendent du cadre le plus général au cadre local
noethérien, où l'énoncé devient reconnaissable.

\pagerange{9}{9}
\subsection*{Reconnaître un isomorphisme sur la fibre spéciale (lemme 3.2, page 9)}

Le feuillet s'ouvre au milieu des hypothèses. Voici la situation telle qu'elle
se reconstitue.\footnote{Le début de l'énoncé est sur sa page 23, qui n'est pas
dans le dossier. Les hypothèses données ici sont celles que la fin de la phrase
et la démonstration exigent ; elles ne sont pas toutes lisibles sur le
feuillet.} $A$ est un anneau préadmissible, $I = I_0$ un idéal de définition,
$(I_i)$ un système fondamental d'idéaux ouverts contenus dans $I_0$, et
$A_i = A/I_i$. $B$ et $C$ sont des $A$-algèbres topologiques dont les topologies
sont définies par les puissances d'idéaux contenant respectivement $IB$ et
$IC$, et l'on suppose $I_iB$ et $I_iC$ fermés — condition automatiquement
vérifiée si $B$ et $C$ sont noethériens.\footnote{Il porte en marge de ce
feuillet que cette hypothèse est inutile, qu'elle « résulte du fait que $IB$
(ou $IC$) [est topologiquement] nilpotent » ; l'abréviation est lue « top. »
sans certitude. L'hypothèse est pourtant conservée ici, et pour une raison
qu'on lit dans sa propre démonstration : c'est elle, et elle seule, qui rend
bijectives les applications canoniques $B \to \varprojlim B_i$ et
$C \to \varprojlim C_i$, dont l'implication (ii) $\Rightarrow$ (i bis) se
sert.} Enfin $u : B \to C$ est un homomorphisme continu de $A$-algèbres, et
$u_i : B_i \to C_i$ l'homomorphisme déduit par passage au quotient.

On suppose une fois pour toutes que $u_0 : B_0 \to C_0$ est un homéomorphisme.
Cette clause figure sur le feuillet à l'intérieur de chacune des conditions ;
il a entouré à l'encre ses quatre occurrences et porté devant chacune une
flèche montante, c'est-à-dire l'ordre de la détacher de la liste. C'est ce qui
est fait ici, et l'énoncé y gagne sa vraie forme.\footnote{Sans cette
correction l'énoncé est redondant : la même clause figurerait dans les quatre
conditions d'une équivalence, ce qui ne peut être que le signe d'une hypothèse
mal placée. La correction est de sa main et elle est juste.}

\emph{Lemme 3.2.} Sous ces hypothèses, les conditions suivantes sont
équivalentes.

\begin{enumerate}
  \item[(i)] $u$ est bijectif.
  \item[(i bis)] $u$ est un isomorphisme d'algèbres topologiques.
  \item[(ii)] Pour tout $i$, $u_i : B_i \to C_i$ est bijectif.
  \item[(ii bis)] Pour tout $i$, en munissant $B_i$ et $C_i$ des filtrations
  $I$-adiques, l'homomorphisme $\mathrm{gr}_I(B_i) \to \mathrm{gr}_I(C_i)$
  induit par $u$ est bijectif.
\end{enumerate}

Si de plus les $C_i$ sont plats sur les $A_i$, ces conditions équivalent
également à

\begin{enumerate}
  \item[(iii)] $u_0 : B_0 \to C_0$ est bijectif.
\end{enumerate}

L'équivalence de (i) et (i bis) est formelle : les topologies de $B$ et $C$
étant définies par les puissances d'idéaux contenant $IB$ et $IC$, une
bijection continue compatible est automatiquement bicontinue. L'implication
(i bis) $\Rightarrow$ (ii) est immédiate, et (ii) $\Rightarrow$ (i bis) résulte
de la bijectivité des applications $B \to \varprojlim B_i$ et
$C \to \varprojlim C_i$, elle-même conséquence du fait que les $I_iB$ et $I_iC$
sont fermés.\footnote{Il renvoie ici à une « rectification à $0_{7.2.4}$ »,
qui n'est pas dans le dossier.} L'équivalence de (ii) et (ii bis) tient à ce
que $I$, étant un idéal de définition, est nilpotent dans chaque $A_i$ : les
filtrations induites sur $B_i$ et $C_i$ sont donc \emph{finies}, et un
morphisme filtré dont le gradué est bijectif est bijectif.

Reste la seule implication qui coûte, (iii) $\Rightarrow$ (ii bis), et c'est
elle qui a besoin de la platitude. Il en trace le diagramme de sa main au bas
du feuillet, et c'est le diagramme qui contient la démonstration :

\begin{tikzcd}
\mathrm{gr}_I(B_i) \arrow[r, "\mathrm{gr}(u_i)"] & \mathrm{gr}_I(C_i) \\
B_0 \otimes_{A_0} \mathrm{gr}_I(A_i) \arrow[r, "u_0 \otimes 1"] \arrow[u] & C_0 \otimes_{A_0} \mathrm{gr}_I(A_i) \arrow[u]
\end{tikzcd}

La flèche verticale de droite est bijective : c'est exactement ce que dit la
platitude de $C_i$ sur $A_i$, sous la forme du critère local de platitude, qui
identifie le gradué de $C_i$ au produit tensoriel de sa fibre spéciale par le
gradué de $A_i$. La flèche verticale de gauche est surjective sans aucune
hypothèse, la surjection $B_0 \otimes \mathrm{gr}(A_i) \to \mathrm{gr}(B_i)$
étant définie pour tout $B$. La flèche du bas est bijective par (iii). Le carré
commute, donc le composé de la flèche du bas suivie de celle de droite est
bijectif ; ce composé étant aussi celui de la flèche de gauche suivie de celle
du haut, la flèche de gauche est injective, donc bijective, et celle du haut est
alors bijective à son tour. C'est (ii bis).

\pagerange{10}{10}
\subsection*{Produits, présentation finie, platitude (3.3 à 3.5, page 10)}

Trois énoncés, dont la fonction est annoncée en une phrase : « il nous sera
commode d'avoir un critère permettant de déduire de la propriété P une
propriété de platitude ».

\emph{Lemme 3.3.} Soient $A$ un anneau, $(P_i)_{i \in I}$ une famille de
$A$-modules, $P = \prod_i P_i$ leur produit, et $M$ un $A$-module de
présentation finie. Alors l'homomorphisme canonique
\[
  M \otimes_A P \longrightarrow \prod_i (M \otimes_A P_i)
\]
est un isomorphisme.

La démonstration est celle du feuillet : on écrit $M$ comme conoyau d'un
homomorphisme $L' \to L$ de modules libres \emph{de type fini}, on note que le
foncteur produit est exact, et l'on est ramené au cas $M = A^n$, où l'énoncé
est trivial. Le membre « de type fini » est une addition manuscrite, et ce
n'est pas une précision de confort : c'est elle qui fait marcher la
démonstration.\footnote{Pour $L = A^{(J)}$ avec $J$ infini, l'application
$A^{(J)} \otimes P \to \prod_i (A^{(J)} \otimes P_i)$ est l'inclusion de
$P^{(J)}$ dans $(\prod_i P_i^{(J)})$, qui n'est pas surjective. Sans « de type
fini » la réduction est donc fausse, et l'hypothèse de présentation finie sur
$M$ perdrait tout emploi. La correction manuscrite répare une lacune réelle de
la version dactylographiée.}

\emph{Corollaire 3.4.} Supposons que tout $A$-module de type fini soit de
présentation finie — par exemple $A$ noethérien. Alors tout produit de
$A$-modules plats est plat.

Il suffit en effet de voir que $P \otimes_A -$ transforme en monomorphisme tout
monomorphisme $M \to N$ de modules \emph{de type fini}, la platitude se testant
sur ces seuls monomorphismes ; le lemme 3.3 ramène alors la question à un
produit de monomorphismes, qui est un monomorphisme.

\emph{Corollaire 3.5.} Supposons $A$ comme en 3.4, par exemple noethérien. Soit
$P$ un $A$-module isomorphe à la limite projective d'un système projectif
\emph{strict}, admettant un système cofinal dénombrable, de $A$-modules
projectifs. Alors $P$ est un $A$-module plat.

Le feuillet dit seulement « on voit facilement que sous les conditions dites,
$P$ est isomorphe à un produit de $A$-modules projectifs, donc plats, d'où le
résultat ». L'argument mérite d'être écrit, car il est court et il est le
ressort de tout le numéro. Le système cofinal dénombrable permet de se ramener
à un système $(P_n)_{n \in \mathbf{N}}$ dont les applications de transition
$P_{n+1} \to P_n$ sont surjectives — c'est ce que veut dire \emph{strict}. Pour
chaque $n$, la suite exacte
\[
  0 \longrightarrow K_n \longrightarrow P_{n+1} \longrightarrow P_n
  \longrightarrow 0
\]
est scindée, puisque $P_n$ est projectif ; donc $P_{n+1} \simeq P_n \oplus K_n$,
et $K_n$, facteur direct d'un projectif, est projectif. En itérant,
$\varprojlim_n P_n \simeq P_0 \times \prod_{n \geqslant 0} K_n$, produit de
modules projectifs. Le corollaire 3.4 conclut.

\pagerange{11}{12}
\subsection*{Le théorème central (3.6, pages 11 et 12)}

\emph{Théorème 3.6.} Soient $A$ un anneau préadmissible muni d'un idéal de
définition $I = I_0$ et d'un système fondamental $(I_i)$ d'idéaux de définition
contenus dans $I_0$, et $B$ une $A$-algèbre topologique. On suppose les
$A_i = A/I_i$ noethériens, et sur $B$ les conditions suivantes :

\begin{enumerate}
  \item[a)] $B$, comme $A$-module topologique, satisfait la condition P ;
  \item[b)] $B_0$ est formellement lisse sur $A_0$ ;
  \item[c)] $B$ est adique ;
  \item[d)] les idéaux $I_iB$ sont fermés dans $B$ — automatiquement vérifié
  si $B$ est noethérien.
\end{enumerate}

Alors :

\begin{enumerate}
  \item[(i)] $B$ est formellement lisse sur $A$.
  \item[(ii)] Soit $C$ une $A$-algèbre topologique satisfaisant à c) et d),
  mais pas nécessairement à a), et soit $u_0 : B_0 \to C_0$ un isomorphisme
  d'algèbres topologiques. Alors il existe un $A$-homomorphisme continu
  $u : B \to C$ se réduisant suivant $u_0$.
\end{enumerate}

Ce dernier étant choisi, les conditions suivantes sont équivalentes :

\begin{enumerate}
  \item[1)] $C$ est formellement lisse sur $A$ ;
  \item[2)] $C$ satisfait la condition a), c'est-à-dire la condition P ;
  \item[3)] les $C_i$ sont plats sur les $A_i$ ;
  \item[4)] $u$ est un isomorphisme de $A$-algèbres topologiques ;
  \item[5)] $C$ est isomorphe à $B$ comme $A$-algèbre topologique.
\end{enumerate}

La liste des cinq conditions est écrite sur le feuillet à propos de $B$ et non
de $C$. C'est une faute de frappe répétée, et la corriger n'est pas une
liberté : à propos de $B$, les conditions 1), 2) et 3) sont des conséquences
immédiates des hypothèses du théorème et l'équivalence est
vide.\footnote{Conditions 1) et 2) à propos de $B$ redisent respectivement
l'assertion (i) et l'hypothèse a) ; la condition 3) résulte de a) par 3.5.
Surtout, sa propre démonstration tranche : elle commence par « supposons 1)
vérifié » et relève alors $u_0^{-1} : C_0 \to B_0$ en un homomorphisme
$v : C \to B$, ce qui exige la lissité formelle de \emph{$C$}, non celle de
$B$. Les cinq conditions portent donc sur $C$, et sont écrites ainsi ici. Les
étiquettes 4) et 5), qui parlent de $u$ et de $C$, sont d'ailleurs les seules
que le feuillet écrive correctement.}

\emph{Démonstration.} L'assertion (i) résulte des hypothèses a) et b), par
l'énoncé 3.1 du numéro, qui ramène la lissité formelle de $B$ sur $A$ à celle
de $B_0$ sur $A_0$.\footnote{Le feuillet écrit « (i) résulte déjà de a) », sans
mentionner b). C'est que la clause b) est une insertion manuscrite : la liste
des hypothèses est renumérotée à la main, un ancien a) est biffé, la clause
« $B_0$ est formellement simple sur $A_0$ » est intercalée en entier à l'encre,
et l'ancien b) devient c). La démonstration, tapée avant l'insertion, n'a pas
été mise à jour. Elle a besoin de b) : la condition P donne la platitude, non
la lissité formelle.} L'existence de $u$ dans (ii) résulte de (i) et de
l'énoncé 1.7, qui n'est pas dans le dossier.

Pour les équivalences : 2) $\Rightarrow$ 1) par (i) appliqué à $C$, qui
satisfait alors a), c) et d), et dont la fibre spéciale $C_0 \simeq B_0$ est
formellement lisse sur $A_0$ ; 4) $\Rightarrow$ 2) et 5) $\Rightarrow$ 2)
trivialement, puisque $B$ satisfait P ; 2) $\Rightarrow$ 3) par 3.5, les $A_i$
étant noethériens. Restent 1) $\Rightarrow$ 4) et 3) $\Rightarrow$ 4).

Supposons 1). Par 1.7, $u_0^{-1} : C_0 \to B_0$ se prolonge en un homomorphisme
continu $v : C \to B$, grâce aux hypothèses c) et d). Or les $B_i$ sont plats
sur les $A_i$, par a) et 3.5. Le lemme 3.2, dans ses conditions (i) et (iii),
les rôles de $B$ et $C$ étant échangés, garantit alors que $v$ est un
isomorphisme d'algèbres topologiques ; $u$ l'est donc aussi. Supposons 3) :
le même argument s'applique dans l'autre sens, la platitude requise par 3.2
(iii) étant cette fois l'hypothèse elle-même. Enfin 4) $\Rightarrow$ 5) est
immédiat.\footnote{Les implications 4) $\Rightarrow$ 5) et 5) $\Rightarrow$ 2)
sont des ajouts manuscrits qui ferment le cycle, le tapuscrit s'arrêtant à
3) $\Rightarrow$ 4). Il note lui-même en clair, à la ligne suivante : « On peut
certainement améliorer la rédaction de la démonstration, évitant certaines
redites. »}

Il annonce aussitôt la restriction qui gouvernera tout le reste : « pour
simplifier, nous nous restreindrons par la suite à des hypothèses
noethériennes », ce qui permet en particulier de se débarrasser de la
« désagréable hypothèse d) ».

\pagerange{12}{13}
\subsection*{Les algèbres de Cohen, et l'unicité (3.7 et 3.8, pages 12 et 13)}

\emph{Définition 3.7.} Soit $A \to B$ un homomorphisme continu d'anneaux
noethériens préadiques. On dit que $B$ est une \emph{algèbre (topologique) de
Cohen} sur $A$ si a) $B$ est séparé et complet, c'est-à-dire un anneau adique ;
b) $B$ satisfait la condition P ; c) $B$ est formellement lisse sur $A$, ou, ce
qui revient au même par 3.1, $B_0 = B/IB$ est formellement lisse sur
$A_0 = A/I$ pour un idéal de définition $I$ de $A$.\footnote{Le feuillet écrit
« $B_0 = B/ID$ » ; il faut lire $B/IB$. La phrase manuscrite qui ferme la
définition court sur deux lignes en interligne et n'est lisible qu'en partie ;
elle porte en tête « $= A/I$ ». — Sur le nom : « algèbre de Cohen » est ici une
notion \emph{relative}, attachée à un homomorphisme, et non la notion absolue
d'\emph{anneau de Cohen} qui a survécu dans la littérature (un anneau local
complet dont l'idéal maximal est engendré par $p$). Rien ici n'est à
confondre avec elle ; et lui-même, quatre feuillets plus loin, propose
d'abandonner le mot.} Il note en marge qu'il conviendrait d'ajouter dès après
3.7 que la notion est stable par extension de la base via produits tensoriels
complétés, « c'est une trivialité ».

\emph{Corollaire 3.8.} Soient $A$ un anneau noethérien, $I$ un idéal, $A$ muni
de la topologie $I$-préadique, $A_0 = A/I$, et $B_0$ une $A_0$-algèbre
topologique. Supposons qu'il existe une algèbre de Cohen $B$ sur $A$ telle que
$B/IB$ soit isomorphe à $B_0$ comme $A_0$-algèbre topologique — ce qui entraîne
que $B_0$ est elle-même de Cohen sur $A_0$. Alors $B$ est unique à isomorphisme
près, compatible avec les isomorphismes donnés des algèbres réduites modulo
$I$ avec $B_0$. De plus, si $B$ est une $A$-algèbre topologique adique telle que
$B/IB \simeq B_0$, les conditions suivantes sur $B$ sont équivalentes :

\begin{enumerate}
  \item[0)] $B$ est une $A$-algèbre de Cohen ;
  \item[1)] $B$ est formellement lisse sur $A$ ;
  \item[2)] $B$ satisfait la condition P ;
  \item[3)] $B$ est plat sur $A$ ;
  \item[3 bis)] pour tout entier $n \geqslant 0$, $B_n = B/I^{n+1}B$ est plat
  sur $A_n = A/I^{n+1}$.
\end{enumerate}

L'unicité est 3.6 (ii) lu à l'envers : si $B$ et $B'$ sont deux algèbres de
Cohen de même fibre spéciale, l'isomorphisme des fibres se relève par (i), et
la condition 4) de la liste dit que le relèvement est un isomorphisme. Des cinq
équivalences, seule celle de 3) et de 3 bis) demande un argument, et il dit où
il est : c'est le critère local de platitude, moyennant les hypothèses
noethériennes.\footnote{Il l'appelle « la platitude supérieure », nom du numéro
des \emph{Éléments} où le critère est démontré. L'implication de 3) vers 3 bis)
est un changement de base ; la réciproque est le critère proprement dit, et
c'est elle qui utilise la complétion.}

Le corollaire 3.8 est le tournant du numéro, et il le souligne. Il montre
« l'intérêt qu'il y a, $A$, $I$ et $B_0$ étant donnés, de disposer d'un
\emph{théorème d'existence} » : l'unicité étant acquise, tout se joue
désormais sur l'existence.

\pagerange{13}{15}
\subsection*{Le problème d'existence, et le critère local (3.9, pages 13 à 15)}

La réduction est classique et il la fait en trois lignes : la construction de
proche en proche, remplaçant $A$ par les $A_i$ et $I$ par ses puissances,
ramène au cas où $I$ est de carré nul, donc où $A$ est discret. « On obtient
alors un problème de nature essentiellement infinitésimale, dont nous ignorons
s'il existe toujours une solution, faute de savoir caractériser la structure des
algèbres de Cohen $B_0$ sur l'anneau discret $A_0$. » C'est l'aveu central du
numéro : le problème général d'existence reste ouvert, et ce qui suit n'en
traite qu'un cas particulier.\footnote{Un premier « Théorème 3.9 » est ici
biffé de deux grands traits diagonaux sur six lignes. Il portait sur $A$ local
noethérien de corps résiduel $k$ et $B_0$ une $k$-algèbre locale noethérienne
complète de corps résiduel $K$ de multiplicité radicielle finie sur $k$,
extension séparable de $k$, et s'arrêtait au milieu d'un mot. Il est remplacé
par le corollaire qui suit, puis le théorème d'existence reparaît à 3.10.}

\emph{Corollaire 3.9.} Soit $A \to B$ un homomorphisme local d'anneaux locaux
noethériens dont l'extension résiduelle est de multiplicité radicielle finie —
par exemple séparable. Les conditions suivantes sont équivalentes, $A$ et $B$
étant munis de leurs topologies habituelles :

\begin{enumerate}
  \item[1)] $B$ est une algèbre de Cohen sur $A$ ;
  \item[2)] $B$ est plat sur $A$, complet, et $B_0 = B/\mathfrak{m}_A B$ est
  une algèbre de Cohen sur $k = A/\mathfrak{m}_A$, c'est-à-dire formellement
  lisse sur $k$.
\end{enumerate}

Une troisième condition est écrite à l'encre dans l'énoncé, « $B$ est
formellement lisse sur $A$ », puis en partie annulée d'un griffonnage ; c'est
celle que le corollaire 3.11 ajoutera, à deux feuillets de là, sous la forme
« $B$ est complet et formellement lisse sur $A$ ». Elle n'est pas retenue ici,
puisqu'il l'a lui-même déplacée.\footnote{L'énoncé dactylographié était une
phrase du type « pour que $B$ soit une algèbre de Cohen, il faut et il suffit
que \ldots » ; la main y plante trois numéros 1), 2), 3) et biffe « il faut et
suffit que », retournant la phrase en liste de conditions équivalentes. La
forme donnée ici est celle que ces corrections produisent.}

Il note entre parenthèses ce qui allège tout : sur un corps, une algèbre
topologique est de Cohen si et seulement si elle est adique et formellement
lisse. La condition P disparaît, puisque tout module sur un corps est libre et
que les quotients $B_0/J_n$ forment un système projectif strict dénombrable.

\emph{Démonstration de} 2) $\Rightarrow$ 1). Tout revient à montrer que $B$
satisfait la condition P, ce qui ramène au cas où $A$ est artinien. Par les
résultats du numéro 2, $B_0$ est régulier ; on choisit un système régulier de
paramètres de $B_0$ — un système de paramètres formant une $B_0$-suite
suffirait — et on le relève en un système $x = (x_1, \ldots, x_r)$ de
paramètres de $B$. Les idéaux $J_n$ engendrés par les $x^n = (x_1^n, \ldots,
x_r^n)$ forment un système fondamental de voisinages de $0$ dans $B$. Les
$B/J_n$ sont alors plats sur $A$, par division successive par un élément
régulier ; et $A$ étant artinien local, un module plat y est
projectif.\footnote{Sur un anneau local artinien, plat entraîne libre, donc
projectif ; le feuillet écrit seulement « comme $A$ est artinien, ils sont donc
projectifs ». — Une phrase manuscrite traverse la page à cet endroit, sur deux
lignes, et n'a pas pu être lue : on n'en relève que les fragments « que \ldots\
3) $\Rightarrow$ 2) », « \ldots\ suffisante » et « \ldots\ d'ailleurs ci-dessous
3.11. ». Rien n'en est reconstruit ici, et l'argument donné ci-dessus est
celui du corps dactylographié seul.}

Suit un NB, barré de six longs traits diagonaux au crayon bleu mais resté
lisible d'un bout à l'autre, et c'est le passage le plus parlant du tapuscrit.
Il y dit d'abord à quoi sert l'hypothèse qu'il vient de faire : « on n'a
utilisé l'hypothèse sur la multiplicité radicielle que pour savoir que $B_0$
est Cohen-Macaulay. L'hypothèse faite est très probablement inutile. » Puis :
« Serre vient de me dire que Mac-Lane a prouvé la réciproque de 2.3. sans
hypothèse de multiplicité rad finie, dans “Note on a relative structure of
$p$-adic fields”, Ann Math t. 41, 1940, page 751-753. Cela vaudrait peut-être
le coup, pour le coup d'œil, que tu regardes cela et donnes le résultat complet
dans 2.3. » En marge du feuillet suivant, à propos de la démonstration de 3.11,
il revient dessus : « Réduction inutile dans le cas “vrai” du
NB. »\footnote{Le NB est barré, non supprimé, et il est resté intégralement
lisible ; il est restitué ici parce qu'il dit à quoi sert une hypothèse que
l'énoncé fait sans la justifier, et qu'aucun autre feuillet ne le dit. La page
suivante porte, à l'encre, l'amorce de la correction qu'il en tire : compléter
l'hypothèse de 3.9 « avec : ou $\mathfrak{m}_B = \mathfrak{m}_A B$ », qui est
exactement l'hypothèse du corollaire 3.12.}

\pagerange{15}{16}
\subsection*{Le théorème d'existence, et le cas non ramifié (3.10 à 3.12, pages 15 et 16)}

\emph{Corollaire 3.10 (existence).} Soient $A$ un anneau local noethérien de
corps résiduel $k$, et $B_0$ une $k$-algèbre qui soit un anneau local
noethérien et une $k$-algèbre de Cohen, c'est-à-dire formellement lisse sur $k$
et complète pour sa topologie ordinaire. Supposons de plus l'extension
résiduelle $K/k$ séparable — ce qui n'est pas une restriction si $B_0 = K$,
mais en est évidemment une en général. Alors il existe une $A$-algèbre
topologique de Cohen $B$ qui se réduit suivant $B_0$ ; elle est unique à
isomorphisme près par 3.8, et c'est un anneau local noethérien complet.

La démonstration est un dévissage en deux cas types, et elle est jolie. Le
théorème de structure de Cohen donne $B_0 \simeq K[[T_1, \ldots,
T_n]]$.\footnote{Le feuillet écrit « en vertu du théorème de structure 3.4 ».
Ce ne peut pas être le 3.4 du numéro, qui porte sur les produits de modules
plats ; c'est le théorème de structure de Cohen, qui appartient au numéro 2. La
référence est donc à lire $2.4$.} Le critère 3.9 2) et la transitivité de la
platitude ramènent alors aux deux cas $1^{\mathrm{o}}$) $B_0 = K$ et
$2^{\mathrm{o}}$) $B_0 = k[[T_1, \ldots, T_n]]$ : dans le premier on invoque le
résultat général de $0_{\mathrm{III}}$ 10.3.1, dans le second il suffit de
prendre $B = \hat{A}[[T_1, \ldots, T_n]]$. Autrement dit : l'extension du corps
résiduel et l'adjonction de variables formelles sont les deux seuls
mouvements, et chacun se relève pour son propre compte.

\emph{Corollaire 3.11.} Sous les conditions de 3.9, l'extension résiduelle
étant de multiplicité radicielle finie, les conditions 1) et 2) équivalent
aussi à : 3) $B$ est complet et formellement lisse sur $A$.

Tout revient à prouver que $B$ est plat sur $A$, et le procédé est une descente.
On fait une extension de la base $A \to A'$, avec $A'$ local, libre de type fini
sur $A$, à extension radicielle convenable ; on trouve alors un $B'$ local sur
$A'$, encore formellement lisse sur $A'$, et cette fois à extension résiduelle
\emph{séparable}. Le théorème d'existence 3.10 s'applique, et 3.8 (équivalence
de 0) et 1)) donne que $B'$ est de Cohen sur $A'$, donc plat sur $A'$. Donc $B$
est plat sur $A$.\footnote{Cette dernière implication est la descente fidèlement
plate de la platitude : $A'$ étant libre de type fini et local sur $A$, il est
fidèlement plat, et $B' = B \otimes_A A'$ plat sur $A'$ entraîne $B$ plat sur
$A$. Le feuillet écrit seulement « Donc $B$ est plat sur $A$ ».}

\emph{Corollaire 3.12.} Soient $A \to B$ un homomorphisme local d'anneaux
locaux noethériens tels que $\mathfrak{m}_B = \mathfrak{m}_A B$, et $K/k$
l'extension résiduelle. Les conditions suivantes sont équivalentes : 1) $B$ est
complet et formellement lisse sur $A$ ; 2) $B$ est complet, plat sur $A$, et
$K/k$ est séparable ; 3) $B$ est une algèbre de Cohen sur $A$. De plus, une
extension séparable $K$ de $k$ étant donnée, il existe une algèbre de Cohen $B$
sur $A$ se réduisant suivant $K$, et $B$ est unique à isomorphisme de
$A$-algèbres près.

C'est le cas d'école, et il vaut la peine de dire ce qu'il est en langage
d'aujourd'hui : l'hypothèse $\mathfrak{m}_B = \mathfrak{m}_A B$ dit que $B$ est
non ramifié sur $A$, la fibre spéciale se réduit alors au seul corps $K$, et la
lissité formelle jointe à la non-ramification est ce qu'on appelle aujourd'hui
la \emph{lissité formelle étale}. L'énoncé dit donc : les extensions locales
complètes non ramifiées de $A$ correspondent exactement aux extensions
séparables de son corps résiduel, et chacune est unique. C'est l'énoncé sur
lequel reposent la construction des anneaux de Cohen et, plus tard, celle de
l'hensélisé strict.

\pagerange{16}{17}
\subsection*{L'énoncé qui reste, et le plan de rédaction abandonné (3.13, pages 16 et 17)}

Avant le dernier corollaire, un paragraphe où il change de registre. « La
rédaction commence à faire un peu touffu, j'espère que tu pourras la rendre
plus lisible. » Puis un regret sur sa propre définition : il n'est « pas
vraiment naturel de supposer $B$ séparée et complète pour sa topologie propre,
mais simplement pour sa topologie $IB$-adique », et il laisse le choix à son
correspondant, en prévoyant l'effet de la modification — dans 3.9 à 3.12 il
faudrait « simplement supprimer à certains endroits le mot “complet” ». Le
feuillet s'achève sur un premier jet manuscrit du corollaire 3.13, aussitôt
annulé de cinq grandes boucles ; le feuillet suivant le reprend au
propre.\footnote{Du premier jet on lit « $A \to B$ \ldots\ local d'anneaux
locaux nœth., \ldots\ de corps \ldots\ à mult.\ deg.\ finie. Alors conditions
équivalentes », et rien de plus : les quatre lignes sont d'une écriture très
rapide et n'ont pas pu être lues davantage. Ce qu'elles apportent n'est donc
pas un contenu mais une chronologie — l'énoncé final a été essayé une fois
avant d'être tapé.}

\emph{Corollaire 3.13.} Soit $A \to B$ un homomorphisme local d'anneaux locaux
noethériens. On suppose l'extension résiduelle $k(B)/k(A)$ de multiplicité
radicielle finie — par exemple séparable, ou de type fini — ou bien
$\mathfrak{m}_B = \mathfrak{m}_A B$.\footnote{Cette dernière clause est une
addition manuscrite, écrite « $k(B) = k(A)B$ ». Lue ainsi elle n'a pas de sens.
La page 15 porte, également à l'encre, l'instruction de compléter l'hypothèse
de 3.9 « avec : ou $r(B) = r(A)B$ » ; c'est la même clause, c'est l'hypothèse
du corollaire 3.12, et c'est la seule lecture qui en fasse une alternative
utilisable. Elle est écrite ici sous cette forme, et la correction est nôtre.}
On munit $A$ et $B$ de leurs topologies habituelles. Alors les conditions
suivantes sont équivalentes :

\begin{enumerate}
  \item[(i)] $B$ est formellement lisse sur $A$ ;
  \item[(ii)] $B$ est plat sur $A$, et $B_0 = B \otimes_A k$ est une algèbre
  universellement régulière sur $k$.
\end{enumerate}

Les deux conditions étant invariantes par remplacement de $B$ par son complété
$\hat{B}$, on se ramène au cas où $B$ est complet, et l'énoncé est alors
contenu dans 3.11, compte tenu de ce que pour $B_0$, formellement lisse et
universellement régulier coïncident — c'est l'énoncé 2.6 du numéro précédent,
qui n'est pas dans le dossier.

C'est l'énoncé sur lequel tout le numéro converge, et c'est celui qu'on
reconnaît aujourd'hui : un homomorphisme local d'anneaux locaux noethériens est
formellement lisse si et seulement s'il est plat à fibre spéciale
géométriquement régulière. C'est la définition même de ce qu'on appelle un
\emph{homomorphisme régulier}, et c'est le point de départ de toute la théorie
qui en découle.

Il le sait, et le dernier paragraphe du dossier est un plan de rédaction contre
lui-même. « Du point de vue pratique, les résultats du présent numéro
consistent essentiellement en 3.13, l'unicité du $B$ correspondant à un $B_0$
donné, et enfin l'existence moyennant la séparabilité de $K/k$. À force de
dévissage et de généralité, cette substance n'apparaît pas assez clairement, et
la terminologie “algèbre de Cohen” introduite concurremment à “formellement
simple” risque plutôt d'embrouiller les idées (d'autant plus que nous
n'utiliserons plus jamais que “formellement simple”). » Suit la recommandation :
énoncer 3.13 en théorème, le démontrer « en deux pages ou trois » en laissant
de côté « tout le fatras anneaux admissibles, algèbres de Cohen etc. »,
commencer par l'existence d'un $B$ satisfaisant (i) et (ii) quand $K/k$ est
séparable, s'y ramener par descente comme dans 3.11, conclure par comparaison
du $B$ type avec le $B$ donné « ce qui donne en même temps l'unicité, et un
point c'est tout ». Une marginale, écrite verticalement sur quatre lignes,
sauve deux choses du naufrage : « garder cependant 3.1 au Nº3. On peut signaler
2.6 pour remarque, à réserver pour l'avenir. »

\pagerange{18}{20}
\section*{III. Questions ouvertes (pages 18 à 20)}

Une feuille nue portant en haut à droite le seul mot « Questions » enveloppe
deux feuillets d'encre. Ce sont des questions qu'il se pose, numérotées 1) à
4), et il faut dire tout de suite dans quel état elles nous parviennent : les
énoncés numérotés ont leur squelette, la prose qui les relie ne se laisse lire
qu'en fragments, et chacun des deux feuillets porte en outre, écrit en oblique
de bas en haut dans toute la marge de gauche, un long bloc d'ajouts postérieurs
— une douzaine de lignes sur le premier, une vingtaine sur le second, avec des
renvois numérotés et des passages soulignés — qui n'ont pas pu être lus du
tout.\footnote{Rien n'est reconstruit ici de ces deux blocs, ni de la prose de
liaison. Ce qui suit est donc, pour cette section seule, moins une lecture
modernisée qu'un relevé de ce qui est sûr ; les questions 3) et 4)
particulièrement sont données dans un état qui ne permet pas de les
interpréter. Sur le premier feuillet on distingue seulement, dans les premières
lignes de la marge, « Faut-il supposer \ldots\ vraiment \ldots ? Peut-être
qu'\ldots », et une formule qui est soit $\Lambda = VA$ soit
$\mathfrak{p} \cap A$.}

\emph{Question 1).} Soit $A$ une algèbre locale noethérienne sur $k$. Est-il
vrai que $A$ est formellement lisse sur $k$ si $A$ est universellement
régulière sur $k$ ?

C'est exactement la réciproque du corollaire 3.13 dépouillée de son hypothèse
de finitude. Sur un corps la platitude est automatique, de sorte que la
condition (ii) de 3.13 se réduit à la régularité universelle de $A$ ; la
question demande donc si l'hypothèse de multiplicité radicielle finie faite sur
l'extension résiduelle peut être levée. Le NB de la page 14 disait déjà qu'il
la croyait superflue.\footnote{Le feuillet porte entre crochets, à la suite de
la question, une glose partiellement illisible qui a la forme d'une définition
— « pour hom. local $A \to B$ d'anneaux locaux noethériens, équivalence entre
“ft simple” et “universellement régulier” » — c'est-à-dire la formulation
générale de 3.13. Le rapprochement fait ci-dessus repose sur cette glose et sur
l'énoncé 3.13 ; il n'est écrit nulle part sur le feuillet.}

\emph{Question 2).} Soit $A$ une algèbre locale noethérienne sur $k$,
formellement lisse sur $k$ pour sa topologie habituelle. Est-il vrai que $A$
est formellement lisse pour sa topologie discrète ?

C'est la question du dossier. Elle est reprise mot pour mot par la première
ligne de la page 21, et les quatre feuillets qui suivent y répondent ; c'est
l'objet de la section IV ci-dessous, et c'est le seul endroit du dossier où
deux campagnes d'écriture se répondent explicitement. Une variante relative est
esquissée à la suite, dont il ne reste que la forme : $A$ plat sur un anneau
$\Lambda$, tous deux noethériens, et la question de savoir si la lissité
formelle des complétés pour les topologies habituelles entraîne celle de $A$
sur $\Lambda$ pour la topologie discrète.

\emph{Question 3).} $B$ plat sur $A$, $A$ et $B$ noethériens. L'ensemble des
points de $\mathrm{Spec} B$ où la fibre est formellement lisse sur le
corps résiduel du point correspondant de $\mathrm{Spec} A$ est-il
ouvert ?

C'est une question d'ouverture du lieu lisse, et c'est tout ce qu'on en peut
dire avec certitude : la moitié des mots qui désignent l'ensemble en question
manque.\footnote{La formulation donnée ici est la seule que les fragments
lisibles — « l'\ldots\ de $\mathrm{Spec} B$ dans l'\ldots\ fermé dans le
$\mathrm{Spec} A$ \ldots\ est ft simple sur le corps résiduel de $A$,
est-il ouvert ? » — rendent plausible, et elle est nôtre. La question se
poursuit sur le feuillet suivant par un crochet ouvert et non refermé.} Deux
choses en revanche sont sûres et valent d'être relevées. La première est la
suite immédiate, sur la page 20 : « faut-il supposer que $A$ et $B$ soient des
“bons anneaux” ». C'est la question qui allait recevoir sa réponse dans le nom
d'\emph{anneau excellent}, et le mot n'est pas encore là. La seconde est une
ligne barrée, réunie par un crochet à l'encre à la mention « Cf IV 2.10.1 » :
c'est le seul renvoi précis à un autre chapitre que portent ces deux
feuillets.\footnote{Le renvoi n'est pas identifié ici : le dossier ne contient
pas le chapitre visé, et la transcription n'a pas ouvert les \emph{Éléments}
imprimés.}

Le feuillet porte encore, entre deux blancs, trois lignes dont on ne relève que
les ancres : une mention « d'extension résiduelle finie » ; une question sur un
« critère jacobien » de Nagata pour une algèbre locale noethérienne ; et une
dernière ligne, qui est celle qu'on retient, où « bons » et « la résolution des
singularités » sont soulignés sur la feuille et où l'on lit que c'est acquis
« en caractéristique $0$ ».\footnote{Le reste de ces trois lignes est illisible,
et notamment ce à quoi se rapporte l'affirmation. On ne peut donc pas dire ce
qu'il tient pour acquis en caractéristique zéro, seulement qu'il l'y tient pour
acquis \emph{pour les bons anneaux} et qu'il souligne, à la suite, la
résolution des singularités. Le dossier n'est pas daté ; la résolution en
caractéristique zéro est le théorème de Hironaka de 1964. Aucune autre
inférence n'est tirée ici de cette ligne.}

\emph{Question 4).} $A \to B$ un homomorphisme d'anneaux noethériens, et de
nouveau une question d'ouverture d'un lieu, portant cette fois sur un quotient.

Elle est tracée en colonne étroite dans la moitié droite du feuillet, et ses
dernières lignes chevauchent l'en-tête imprimé du papier. Elle n'est pas
restituable : ce qui se lit en est une localisation $A' = A_{\mathfrak{p}}$, un
$B'$, un $\Lambda \subset B'$, un produit tensoriel $B \otimes_{A'}$ et la
clause finale « $B'$ plat sur $A'$ », dans un ordre que la page ne permet pas
de reconstituer.\footnote{Les identifications $A'$ et $B'$ sont elles-mêmes
données comme douteuses par la transcription. Rien n'est proposé ici.}

\pagerange{21}{24}
\section*{IV. $K[[t_1,\ldots,t_n]]$ est-il formellement lisse sur $k$ ? (pages 21 à 24)}

Quatre feuillets d'encre, écrits d'un trait, sur le verso de feuilles à en-tête
de sa dactylographe. Ils s'ouvrent sur la question 2) de la page 19, reposée
dans les mêmes termes : une algèbre locale séparée et complète, formellement
lisse pour la topologie habituelle, l'est-elle pour la topologie discrète ?
Trois points d'interrogation ferment la ligne. La réponse tient en deux cas, et
ils sont opposés.

\pagerange{21}{21}
\subsection*{Caractéristique $p > 0$ : les p-bases suffisent (page 21)}

En caractéristique $p$, la lissité formelle d'une extension se lit sur les
p-bases, et c'est ce qu'il utilise. Soit $A \simeq K[[t_1, \ldots, t_n]]$ avec
$K/k$ une extension, et supposons $K$ fini sur $K^p$ — c'est dire que $K$
possède une p-base finie sur $k$. On prend alors une p-base $(x_1, \ldots,
x_m)$ de $K$ et on lui adjoint les variables $t_1, \ldots, t_n$. Le point est
que la complétion selon les $t_i$ ne détruit pas la p-indépendance : le système
$(x_1, \ldots, x_m, t_1, \ldots, t_n)$ est une p-base de $A$ sur $k$, et $A$
est formellement lisse sur $k$, pour la topologie discrète comme pour
l'autre.\footnote{Le feuillet accompagne cette conclusion d'une formule dont la
parenthèse n'est pas refermée, « $k(A^p \simeq K \otimes_{K^p} A^p$ », qui dit
que $K$ et $A^p$ sont linéairement disjoints sur $K^p$ — c'est-à-dire
exactement que les $t_i$ restent p-indépendants au-dessus de $K$. Le terme
après le signe $\simeq$ est repassé et n'a pas été lu avec certitude. La marge
porte, isolément : « $\Omega_{K/k}$ est libre sur $K$ », qui est la traduction
de l'hypothèse de p-base et qui annonce l'invariant dont tout dépendra.} Il
ajoute le renvoi « [cf. “Diff”] », qui vise la chemise « Différentielles »
quatre feuillets plus loin : c'est le dossier qui se cite lui-même.

\pagerange{22}{24}
\subsection*{Caractéristique nulle : « Je suis pessimiste » (pages 22 à 24)}

Le cas contraire s'ouvre sur ces quatre mots, et sur une conjecture qu'il pose
avant de la démontrer : dès que $\dim A > 0$, $A$ n'est \emph{jamais}
formellement lisse sur $k$ pour la topologie discrète.\footnote{Le feuillet
porte « Je présume que de dim $A \geqslant 0$ », lecture dont la moitié est
illisible. La dimension nulle doit être exclue : $A$ se réduit alors au corps
$K$, et en caractéristique nulle toute extension de corps est séparable, donc
formellement lisse pour la topologie discrète. C'est $\dim A > 0$ qui est écrit
ici, et la correction est nôtre.}

Le mécanisme est celui-ci, et il est entièrement contenu dans une question de
place. Si $A$ est formellement lisse sur $k$ pour la topologie discrète, la
suite exacte fondamentale des différentielles
\[
  \Omega^1_{k} \otimes_{k} A \longrightarrow \Omega^1_{A}
  \longrightarrow \Omega^1_{A/k} \longrightarrow 0
\]
se complète à gauche par un zéro et se scinde ; en particulier il existe une
rétraction $\Omega^1_A \to \Omega^1_k \otimes_k A$, et en la composant avec la
différentielle de $A$ on obtient une dérivation
\[
  D : A \longrightarrow \Omega^1_{k} \otimes_{k} A
\]
qui induit $d_{k}$ sur $k$. C'est cette dérivation que les pages 22 à 24
cherchent, et c'est son inexistence qui est l'obstruction.

\emph{L'exemple.} Prenons $k = \mathbf{Q}(x_i : i \in I)$ avec $I$ infini, et
$A = k[[t]]$. Alors $\Omega^1_{k}$ est un $k$-espace vectoriel de base
$(dx_i)_{i \in I}$, de sorte que
\[
  \Omega^1_{k} \otimes_{k} A \simeq A^{(I)} \subsetneq \prod_{i \in I} A\,dx_i ,
\]
la somme directe étant \emph{strictement} contenue dans le produit dès que $I$
est infini. Toute la démonstration est là.

\emph{L'argument.} Supposons donnée une telle $D$, et supposons d'abord
$D(t) = 0$. Alors pour tout $n$, $D(\mathfrak{m}^n) \subset \mathfrak{m}^n
\cdot (\Omega^1_k \otimes_k A)$, puisque $\mathfrak{m}^n = t^n A$ et que
$D(a t^n) = D(a)\,t^n$. La dérivation $D$ est donc automatiquement continue
pour les topologies $\mathfrak{m}$-adiques ; le but étant
$\mathfrak{m}$-adiquement séparé, $D$ est entièrement déterminée par ses
valeurs sur les polynômes, et l'on a nécessairement
\[
  D\Big(\sum_{\lambda} a_{\lambda} t^{\lambda}\Big)
  \;=\; \sum_{\lambda} d_{k}(a_{\lambda})\, t^{\lambda} .
\]
Or le membre de droite n'a aucune raison d'appartenir à $A^{(I)}$. Il suffit,
pour le prendre en défaut, de choisir une suite d'indices distincts $i_0, i_1,
i_2, \ldots$ et des coefficients $a_n \in k$ tels que $d_k(a_n)$ ait une
composante non nulle sur $dx_{i_n}$ — par exemple $a_n = x_{i_n}$ — et de poser
\[
  f = \sum_{n \geqslant 0} a_n t^{n} .
\]
La suite des sommes partielles de $D(f)$ est de Cauchy pour la topologie
$\mathfrak{m}$-adique et converge donc dans le complété de $A^{(I)}$, qui
contient $A^{(I)}$ strictement ; mais sa limite a une infinité de composantes
non nulles et n'appartient pas à $A^{(I)}$. La dérivation $D$ ne peut donc pas
exister, et $A$ n'est pas formellement lisse sur $k$ pour la topologie
discrète.\footnote{Le feuillet 24 s'arrête sur la construction de $f = \sum a_n
t^n$ ; la conclusion n'est pas écrite. Elle est fournie ici, et la rédaction
de l'argument de continuité l'est aussi : les pages 22 et 23 disent seulement
que $D(a)$ est « une combinaison linéaire finie de certains $\ldots x_i$,
$i \in J$ », ce qui est la même observation prise par l'autre bout. — La
réduction au cas $D(t) = 0$ est ce que les bas de pages 22 et 23 mènent, en
passant à un sous-corps $k_0 = \mathbf{Q}(x_i : i \in J)$ engendré par les
indices $J$, en nombre fini, qui apparaissent dans $D(t)$ ; mais la prose qui
l'exécute est un palimpseste et n'est pas lisible. Seule sa forme est donnée
ici. Le feuillet 22 porte en outre une prescription « $D(x_i) = 1/\alpha\,x_i$ »
dont la lettre au dénominateur est une seule lettre, lue $\alpha$ sans
certitude ; elle n'est pas utilisée ici.}

\emph{Ce que la page 23 énonce.} Entre la conjecture et le contre-exemple, il
pose une proposition dont l'hypothèse est nette et dont la conclusion ne l'est
pas : $k$ un corps, $K$ une extension telle que $\Omega_{K/k}$ soit \emph{de
dimension finie} sur $K$, et $A = K[[t_1, \ldots, t_n]]$. La conclusion est
que $A$ est alors formellement lisse sur $k$ pour la topologie
discrète.\footnote{La moitié droite de l'énoncé est illisible, et la conclusion
donnée ici est nôtre. Elle est celle que l'hypothèse force : si $\Omega_{K/k}$
est de dimension finie sur $K$, alors $\Omega_{K/k} \otimes_K A$ est un
$A$-module libre de \emph{rang fini}, somme directe et produit coïncident, et
l'obstruction exhibée trois lignes plus bas s'évanouit. Elle est aussi celle
que la suite du feuillet suppose, puisqu'il y cherche une $D$ à valeurs dans
$\Omega_{K/k} \otimes_K A$ induisant $d_{K/k}$ et qu'il conclut sur un
« l'hypothèse ! ». Aucune autre conclusion ne rendrait compte de ces deux
choses.}

C'est ce qui unifie les deux cas, et c'est le vrai contenu de ces quatre
feuillets : l'invariant qui décide n'est ni la caractéristique ni la dimension,
c'est \emph{la dimension de $\Omega_{K/k}$ sur $K$}. En caractéristique $p$,
l'existence d'une p-base finie la rend finie, et tout va bien. En
caractéristique nulle sur un corps engendré par une infinité de variables, elle
est infinie, et rien ne va. Les trois campagnes du dossier — la marginale de la
page 21, la proposition de la page 23, le contre-exemple des pages 22 et 24 —
disent la même chose de trois manières, sans que le dossier le dise jamais.

\pagerange{25}{28}
\section*{V. Différentielles : le carré $k \to K$ et son contre-exemple (pages 25 à 28)}

Une feuille crème par ailleurs blanche porte en haut à droite, souligné, le
seul mot « Différentielles » : c'est la chemise à laquelle renvoyait la page
21. Ce qu'elle enveloppe est d'une autre campagne, au crayon, d'aspect
postérieur, et n'occupe que deux feuillets.\footnote{Les pages 27 et 29, versos
de ces deux feuillets, portent un imprimé sur la topologie de l'espace des
répartitions et le groupe des idèles, sans rapport avec le chapitre et sans
annotation de lui. Elles ne sont pas lues.}

\emph{La remarque.} Soit $K$ un corps de caractéristique $p > 0$ et $k$ un
sous-corps. Il considère le carré

\begin{tikzcd}
k' = k \cap K^{p} \arrow[r, no head] \arrow[d, no head] & K^{p} \arrow[d, no head] \\
k \arrow[r, no head] & K
\end{tikzcd}

et pose la question de l'injectivité de l'application canonique
\[
  \Omega^{1}_{k/k \cap K^{p}} \otimes_{k} K \longrightarrow \Omega^{1}_{K} ,
\]
dans le prolongement de la suite $\Omega^{1}_{k'} \otimes_{k'} K \to
\Omega^{1}_{k} \otimes_{k} K \to \Omega^{1}_{K}$.\footnote{Le feuillet dessine
le carré avec $k \to K$ en haut, muni d'une pointe de flèche, et les trois
autres côtés sans pointe ; il est redressé ici de manière que les quatre côtés
soient des inclusions et que le sous-corps commun soit en haut à gauche. Le
contenu est le même et les deux dispositions sont également légitimes. — La
prose qui entoure la question est presque entièrement illisible : on y
distingue des différentielles $d_{k/k'}x$, $d_{K/k'}\gamma$, $d_{K/K^p}x$,
$d_{K/K^p}\gamma$ et, à deux reprises, une norme $N_{K/k}$ dont le rôle ne se
laisse pas reconstituer.}

C'est la question de la validité, pour une extension de corps en
caractéristique $p$, de la première suite exacte fondamentale des
différentielles. Elle est équivalente à une condition de séparabilité, et elle
est fausse en général : il en construit un contre-exemple, et c'est le seul
passage complètement démontré de ces deux feuillets.

\emph{Le contre-exemple.} Soit $\mathbf{F}_0$ le corps premier à $p$ éléments.
On prend $x$, $y$, $a_1, \ldots, a_{p-1}$ algébriquement indépendants sur
$\mathbf{F}_0$, et l'on pose
\[
  a_0 = y - \sum_{i=1}^{p-1} a_i x^{i} ,
  \qquad\text{de sorte que}\qquad
  y = \sum_{i=0}^{p-1} a_i x^{i} .
\]
On pose alors $K = \mathbf{F}_0(x, a_0^{1/p}, a_1^{1/p}, \ldots,
a_{p-1}^{1/p})$, de sorte que $K^p = \mathbf{F}_0(x^p, a_0, a_1, \ldots,
a_{p-1})$ : tous les $a_i$ sont dans $K^p$, tandis que $x$, qui est de degré
$p$ sur $K^p$, n'y est pas. Enfin $k = \mathbf{F}_0(x,y)$, qui est bien
contenu dans $K$ puisque $y = \sum a_i x^i$. On note $k' = k \cap K^p$, et l'on
observe que $x^p \in k'$ et $x \notin k'$, de sorte que $k'(x)$ est de degré
$p$ sur $k'$, de base $1, x, \ldots, x^{p-1}$.

Deux vérifications. La première est que dans $\Omega^1_K$, les différentielles
de $x$ et de $y$ sont \emph{liées} : les $a_i$ étant dans $K^p$, ils sont
annulés par toute dérivation de $K$, et la relation $y = \sum a_i x^i$ donne
\[
  d y = \Big(\sum_{i=1}^{p-1} i\,a_i x^{i-1}\Big)\,d x
  \qquad \text{dans } \Omega^1_K .
\]
La seconde est que dans $\Omega^1_{k/k'}$ elles ne le sont pas, parce que le
couple $(x,y)$ est p-indépendant sur $k'$. C'est l'argument de la page 28, et
il est par l'absurde. Si $(x,y)$ n'était pas p-indépendant sur $k'$, alors
$k = k'(x,y)$ serait de degré au plus $p$ sur $k'$, donc égal à $k'(x)$, et
$y$ s'écrirait $y = \sum_{i=0}^{p-1} a'_i x^{i}$ à coefficients $a'_i \in k'$.
Mais $x \notin K^p$, donc $1, x, \ldots, x^{p-1}$ est aussi une base de
$K^p(x)$ sur $K^p$, et les deux écritures de $y$ à coefficients dans $K^p$
coïncident : $a'_i = a_i$ pour tout $i$. On aurait donc $a_i \in k =
\mathbf{F}_0(x,y)$, ce qui est absurde, les $a_i$ ayant été pris
algébriquement indépendants de $x$ et de $y$.

Donc $dx$ et $dy$ engendrent dans $\Omega^1_{k/k'} \otimes_k K$ un sous-espace
de dimension deux sur $K$, dont l'image dans $\Omega^1_K$ est la droite
$K\,dx$ : l'application n'est pas injective, et la remarque est
tranchée.\footnote{Les deux feuillets portent, autour de cet argument, plusieurs
essais biffés du même exemple — une version en $\alpha$, $y - \alpha x$ et
$b = a^{1/p}$, et une réduction où l'on pose $a_i = 0$ pour
$2 \leqslant i \leqslant p-1$, qui ramène au cas $y = a_0 + a_1 x$ — ainsi
qu'une remarque finale
indiquant que $\mathbf{F}$ peut être un corps fini. Ils ne sont pas rendus
ici : ce sont des variantes du même contre-exemple, et aucune n'est écrite
assez complètement pour être vérifiée. La reconstitution donnée ci-dessus est
celle des fragments cohérents des pages 26 et 28 ; l'identification du corps de
base $\mathbf{F}$, que le feuillet écrit tantôt $\mathbf{F}$ tantôt
$\mathbf{F}_0$, est nôtre.}

\end{document}
